\begin{table}[h!]
\centering
\caption{\textbf{Treated Sample Construction and Coverage}}
\label{tab:SampleConstruction}
\resizebox{\textwidth}{!}{%
\begin{tabular}{lcccc}
\toprule
& \multicolumn{2}{c}{Down-Regulation (III to II)}
& \multicolumn{2}{c}{Deregulation (II to I)} \\
\cmidrule(lr){2-3} \cmidrule(lr){4-5}
& $N$ & \% & $N$ & \% \\
\midrule
\multicolumn{5}{l}{\textbf{A. Construction of treated sample}} \\[2pt]
Federal Register reclassified device types
  & 31 & 100 & 356 & 100 \\
Excluding atypical reclassifications\textsuperscript{a}
  & 29 & 94 & n/a & n/a \\
FDA-initiated only (excluding petition-driven cases)
  & 12 & 39 & n/a & n/a \\
With $\geq$1 pre-reclassification patent
  & 12 & 39 & 219 & 62 \\
Final treated device-type sample
  & 12 & 39 & 219 & 62 \\
\\[-2pt]
\multicolumn{5}{l}{\textbf{B. Estimation-sample coverage by outcome}} \\[2pt]
Patent counts, forward citations, safety emphasis
  & 12 & 100 & 219 & 100 \\
Patent market value\textsuperscript{b}
  & 12 & 100 & 219 & 100 \\
Device submissions\textsuperscript{c}
  & 12 & 100 & n/a & n/a \\
Adverse events
  & 12 & 100 & 219 & 100 \\
Price and utilization\textsuperscript{d}
  & 12 & 100 & 96 & 44 \\
Public-firm heterogeneity\textsuperscript{e}
  & 12 & 100 & 219 & 100 \\
\bottomrule
\end{tabular}}

\vspace{3mm}

\caption*{\footnotesize
\textit{Note:} Panel~A traces the treated device-type sample from the Federal Register reclassification universe to the final analysis sample. Panel~B reports, for each outcome group, the number and share of treated device types that enter the corresponding estimating sample. Patent counts, forward citations, and safety-emphasis measures are constructed from all USPTO grants in the Dimensions patent database, covering both public and private firms. In device-type-years with no patenting, these patent outcomes are coded as zero. In device-type-years with no reported adverse events, adverse-event counts are coded as zero. All such zero-valued observations enter the regressions. Tables~\ref{tab:Combined_DropNoCounts} report robustness estimates excluding device types whose outcome sum equals zero over the full sample period; because this restriction is applied separately by outcome, sample sizes can differ slightly across specifications.\\[3pt]
\textsuperscript{a}Two atypical III-to-II cases are excluded: LTT (antimicrobial susceptibility panels), a preamendment-to-Class~II reclassification rather than a III-to-II reclassification; and HRD (soft-lens thermal sterilizer), a commercially obsolete product with no device submissions after 1983, well before its 1997 reclassification.\\
\textsuperscript{b}Patent market values follow \citet{kogan:2017} and are observed only for publicly traded assignees. Device-type-years in which patenting is attributable only to private firms are coded as zero rather than excluded, so the device-type sample is identical across patent-based outcomes.\\
\textsuperscript{c}Class~I devices are generally exempt from premarket submission requirements, so device-submission outcomes are not defined for II-to-I reclassifications.\\
\textsuperscript{d}Requires linkage to MarketScan claims (1996--2013) and outpatient procedure mappings, restricting the sample to post-1996 reclassification events with usable procedure links.\\
\textsuperscript{e}Uses firm assets from CRSP/Compustat and is the only specification restricted to publicly listed firms.}
\end{table}